\documentclass{amsbook}

\usepackage{amssymb}
\usepackage{hyperref}
\usepackage{enumitem}
\usepackage{geometry}                % See geometry.pdf to learn the layout options. There are lots.
\geometry{a4paper}                   % ... or a4paper or a5paper or ... 

\newtheorem{theorem}{Theorem}[chapter]
\newtheorem{lemma}[theorem]{Lemma}

\theoremstyle{definition}
\newtheorem{definition}[theorem]{Definition}
\newtheorem{example}[theorem]{Example}
\newtheorem{xca}[theorem]{Exercise}

\theoremstyle{remark}
\newtheorem{remark}[theorem]{Remark}

\numberwithin{section}{chapter}
\numberwithin{equation}{chapter}

\begin{document}

\frontmatter

\title{Minqi's Solutions to the Problems of Peter J. Bickel and Kjell A. Doksum (2015), Mathematical Statistics: Basic Ideas and Selected Topics, Volume I (2nd Edition)}

\author{Minqi Pan\footnote{Website: http://www.minqi-pan.com}}

\date{\today}

\maketitle

\setcounter{page}{4}

\tableofcontents

\mainmatter

\chapter{Statistical Models, Goals, and Performance Criteria}

\section{Problem 1.1.1: Examples of Statistical Models}
Give a formal statement of the following models identifying the probability laws of the data and the parameter space. State whether the model in question is parametric or nonparametric.
\begin{enumerate}[label=(\alph*)]
\item A geologist measures the diameters of a large number $n$ of pebbles in an old stream bed.
Theoretical considerations lead him to believe that the logarithm of pebble diameter is normally distributed with mean $\mu$ and variance $\sigma$.
He wishes to use his observations to obtain some information about $\mu$ and $\sigma^2$
but has in advance no knowledge of the magnitudes of the two parameters.
\item A measuring instrument is being used to obtain $n$ independent determinations of a physical constant $\mu$.
Suppose that the measuring instrument is known to be biased to the positive side by $0.1$ units.
Assume that the errors are otherwise identically distributed normal random variables with known variance.
\item In part (b) suppose that the amount of bias is positive but unknown. Can you perceive any difficulties in making statements about $\mu$ for this model?
\item The number of eggs laid by an insect follows a Poisson distribution with unknown mean $\lambda$.
Once laid, each egg has an unknown chance $p$ of hatching and the hatching of one egg is independent of the hatching of the others.
An entomologist studies a set of $n$ such insects observing both the number of eggs laid and the number of eggs hatching for each nest.
\end{enumerate}
\begin{proof}[Answer]
\begin{enumerate}[label=(\alph*)]
\item Let sample space $\Omega$ be defined by\[
\Omega\equiv\{(P_1,\dots,P_n):P_i\text{ is a pebble of the old stream}\}.
\]Let $\mathbf{X}=(X_1,\dots,X_n):\Omega\to\mathbb{R}^n$ be the random vector on $\Omega$ defined by\[
\mathbf{X}(P_1,\dots,P_n)\equiv(\log(\text{diameter of }P_i),\dots,\log(\text{diameter of }P_n)).
\]We assume that $\mathbf{X}\sim P_{\mu,\sigma^2}$, where $P_{\mu,\sigma^2}$ is the probability distributions on $\mathbb{R}^n$ according to which $X_1,\dots,X_n$ are independent and identically distributed with a common $\mathcal{N}(\mu,\sigma^2)$ distribution. Let the parameter space $\Theta$ be defined by\[
\Theta\equiv\{(\mu,\sigma^2):\mu\in\mathbb{R},\sigma^2>0\}.
\]Finally, let model $\mathcal{P}$ be defined by\[
\mathcal{P}\equiv\{P_{\mu,\sigma^2}:(\mu,\sigma^2)\in\Theta\}.
\]$\mathcal{P}$ is parametric.
\item Let sample space $\Omega$ be defined by\[
\Omega\equiv\{(D_1,\dots,D_n)\in\mathbb{R}^n:D_i\text{ is a determination of the physical constant}\}.
\]Let $\mathbf{X}=(X_1,\dots,X_n):\Omega\to\mathbb{R}^n$ be the random vector on $\Omega$ defined by\[
\mathbf{X}(D_1,\dots,D_n)\equiv(D_1,\dots,D_n).
\]Let $\mathbf{Y}=(Y_1,\dots,Y_n):\Omega\to\mathbb{R}^n$ be the random vector on $\Omega$ defined by\[
\mathbf{Y}=\mathbf{X}-\mu.
\]We assume that $\mathbf{Y}\sim P_{\sigma^2}$, where $P_{\sigma^2}$ is the probability distributions on $\mathbb{R}^n$ according to which $Y_1,\dots,Y_n$ are independently and identically distributed with a common $\mathcal{N}(+0.1,\sigma^2)$ distribution. Therefore\[
X_i\sim\mathcal{N}(\mu+0.1,\sigma^2),\quad\forall i=1,\dots,n.
\]That is to say, $\mathbf{X}\sim P_{\mu+0.1,\sigma^2}$, where $P_{\mu+0.1,\sigma^2}$ is the probability distributions on $\mathbb{R}^n$ according to which $X_1,\dots,X_n$ are independent and identically distributed with a common $\mathcal{N}(\mu+0.1,\sigma^2)$ distribution. Let the parameter space $\Theta$ be defined by\[
\Theta\equiv\{(\mu+0.1,\sigma^2):\mu+0.1\in\mathbb{R},\sigma^2>0\}.
\]Finally, let model $\mathcal{P}$ be defined by\[
\mathcal{P}\equiv\{P_{\mu+0.1,\sigma^2}:(\mu+0.1,\sigma^2)\in\Theta\}.
\]$\mathcal{P}$ is parametric.
\item Yes. The difficulty is that we need one more parameter $\mu'$, and the parametrization becomes unidentifiable. Specifically, we need to assume that $\mathbf{Y}\sim P_{\mu',\sigma^2}$, where $P_{\mu',\sigma^2}$ is the probability distributions on $\mathbb{R}^n$ according to which $Y_1,\dots,Y_n$ are independently and identically distributed with a common $\mathcal{N}(\mu',\sigma^2)$ distribution. Therefore\[
X_i\sim\mathcal{N}(\mu+\mu',\sigma^2),\quad\forall i=1,\dots,n.
\]That is to say, $\mathbf{X}\sim P_{\mu,\mu',\sigma^2}$, where $P_{\mu,\mu',\sigma^2}$ is the probability distributions on $\mathbb{R}^n$ according to which $X_1,\dots,X_n$ are independent and identically distributed with a common $\mathcal{N}(\mu,\mu',\sigma^2)$ distribution. Let the parameter space $\Theta$ be defined by\[
\Theta\equiv\{(\mu,\mu',\sigma^2):\mu\in\mathbb{R},\mu'\in\mathbb{R},\sigma^2>0\}.
\]Finally, let model $\mathcal{P}$ be defined by\[
\mathcal{P}\equiv\{P_{\mu,\mu',\sigma^2}:(\mu,\mu',\sigma^2)\in\Theta\}.
\]$\mathcal{P}$ is parametric.
\item Let sample space $\Omega$ be defined by\[
\Omega\equiv\{(N_1,\dots,N_n):N_i\text{ is a nest}\}.
\]Let $\mathbf{X}:\Omega\to\mathbb{R}^n$ be the random vector on $\Omega$ defined by\[
\begin{split}
\mathbf{X}&\equiv(X_1,\dots,X_n),\\
X_i(N_1,\dots,N_n)&\equiv\text{the number of eggs laid in }N_i.
\end{split}
\]Let $\mathbf{Y}:\Omega\to\mathbb{R}^n$ be the random vector on $\Omega$ defined by\[
\begin{split}
\mathbf{Y}&\equiv(Y_1,\dots,Y_n),\\
Y_i(N_1,\dots,N_n)&\equiv\text{the number of eggs hatching in }N_i.
\end{split}
\]We assume that $(\mathbf{X},\mathbf{Y})\sim P_{\lambda,p}$, where $P_{\lambda,p}$ is the probability distributions on $\mathbb{R}^{2n}$ according to which $(X_1,Y_1),\dots,(X_n,Y_n)$ are independent and identically distributed with a common distribution function $F_{\lambda,p}$, where\[
F_{\lambda,p}(x,y)=\mathcal{P}(x;\lambda)\cdot\mathcal{B}(y;x,p)
\]Let the parameter space $\Theta$ be defined by\[
\Theta\equiv\{(\lambda,p):\lambda>0,p\in[0,1]\}.
\]Finally, let model $\mathcal{P}$ be defined by\[
\mathcal{P}\equiv\{P_{\lambda,p}:(\lambda,p)\in\Theta\}.
\]$\mathcal{P}$ is parametric.
\end{enumerate}
\end{proof}

\end{document}
